% MY51990C
% Unterlassung der Reanimation
% 14.0
% 2013-11-30
\label{m:reanimation-unterlassung}
\Ca{Unter den folgenden Voraussetzungen darf (oder muss) eine
Reanimation unterlassen
werden:\index{Reanimation!Unterlassung
der}\index{Wiederbelebung|see{Reanimation}}}           
\begin{itemize}
  \item \EE{Verwesung}
  \item \EE{Fäulnis}
  \item \EE{Mumifizierung}
  \item \EE{Skelettierung}
  \item \EE{Verletzungen, die nicht mit dem Leben
  vereinbar sind}
  \item \EE{Tierfraß}\footnote{Unter
  \E{\enquote{Tierfraß}} versteht man den Verzehr eines Kadavers
  durch Kadaverfresser. Ein frischer Hundebiss oder
  Ähnliches fällt \E{nicht} darunter. Maden können
  sich im Gewebe von Lebenden einnisten und sind daher
  auch kein sicheres Todeszeichen.}
  \item (Verbindliche) \EE{Patientenverfügung}\footnote{\E{Patientenverfügung}: Bei Notfallversorgungen darf \EE{nicht nach einer Verfügung gesucht
  werden, wenn dadurch das Leben oder die Gesundheit des Patienten ernstlich
  gefährdet werden würde}.
  \EE{Aber:} Wenn \E{zweifelsfrei} eine verbindliche
  Patientenverfügung vorliegt (z.\,B. weil eine betreuende
  Pflegekraft diese in der Zwischenzeit gefunden und die
  Verbindlichkeit bestätigt hat), muss dieser Folge
  geleistet werden.
Liegt eine \E{beachtliche Patientenverfügung} vor, so ist im
\E{Einzelfall} zu entscheiden: Das Wohl des Patienten bleibt
oberstes Gebot und die Verfügung soll in die
Entscheidung einfließen.}\index{Patientenverfügung!Unterlassung der
  Reanimation} (\FullPrettyRef{AASS-topic:patientenverfuegung}; 
  Bis zur zweifelsfreien Klärung der Situation muss
  jedenfalls eine Reanimation begonnen bzw. fortgeführt
  werden!)
  \item (Abbruch der Reanimationsmaßnahmen aufgrund
  von \EE{Erschöpfung})\footnote{Erschöpfung darf im
  professionellen Umfeld keine Rolle spielen, für eine
  entsprechende Ablösung ist bei Bedarf zu sorgen.}
  \item Unvereinbarkeit mit
  \EE{Selbstschutz}\footnote{Es müssen jedoch alle Maßnahmen
  getroffen werden, die ohne Gefährdung des Selbstschutzes
  möglich sind. So wird zum Beispiel für Ersthelfer
  empfohlen, wenn
  eine Mund-zu-Mund-Beatmung nicht zumutbar ist, zumindest 
  eine Herzdruckmassage durchzuführen.}
\end{itemize}
